\subsection{Deterministic likelihood replacement for every fixed $K$}

\begin{lemma}[Iterative likelihood replacement certifies a weak start]
\label{lem:iterative-swap-certification}
Let $K\ge2$ be fixed, put $P_{\max}=\max_{a,b}P_{ab}$, and suppose
$nP_{\max}\to\infty$. Assume, for constants $0<c\le C<\infty$,
$\pi_{\min}>0$, and $s>0$, that
\begin{align}
 n_a&\ge\pi_{\min}n,\qquad
 cP_{\max}\le q_{ij},p_{aj}\le CP_{\max},
 \label{eq:swap-profile-box}\\
 \sum_i\|q_i-p_{z_i,*}\|_1&\le a_n n^2P_{\max},
 \qquad a_n\longrightarrow0,
 \label{eq:swap-average-error}\\
 \min_{a\ne b}\|p_{a*}-p_{b*}\|_1&\ge s nP_{\max}.
 \label{eq:swap-profile-separation}
\end{align}
Start at any $K$ fitted profiles in the convex hull of the rows of $q$
and any feasible simplex weights. A fixed weight floor
$f\in[0,1/K)$ is allowed. Define
\begin{align}
 D(x\Vert y)&=\sum_j\{x_j\log(x_j/y_j)-x_j+y_j\},\notag\\
 J(\widehat p)&=\sum_i\min_{k\le K}D(q_i\Vert\widehat p_{k*}),
 \label{eq:swap-hard-cost}\\
 F(\widehat\pi,\widehat p)
 &=-\sum_i\log\sum_{k=1}^K\widehat\pi_k
       e^{-D(q_i\Vert\widehat p_{k*})},
 \label{eq:swap-soft-cost}\\
 T_n&=\sum_{i,j}q_{ij},\qquad
 \zeta_n=(\log n)^{-1/2},\qquad
 M_n=\lceil(\log n)^2\rceil.
 \label{eq:swap-stopping-constants}
\end{align}

At each replacement round form the following $K+1$ proposals.
One proposal starts at the current profiles with weights $1/K$.
For each component $r$, retain the other $K-1$ profiles and choose
\begin{equation}
 c_r\in\argmax_{c\in[n]}
 \sum_i\left(\ell_i(q_c)-\max_{k\ne r}
                         \ell_i(\widehat p_{k*})\right)_+.
 \label{eq:swap-candidate-rule}
\end{equation}
Replace profile $r$ by $q_{c_r}$ and initialize all $K$ weights at
$1/K$. Fit each proposal by any prescribed finite number of exact EM
updates, including zero updates, whose M-steps maximize over the stated
weight simplex. A zero-mass profile retains its previous value.
Select the endpoint with largest $\mathcal L$ and accept it only when
its objective increase is at least $\zeta_nT_n$; otherwise stop.
After $M_n$ accepted rounds, return the current fit without further
replacement. Candidate scans in \eqref{eq:swap-candidate-rule} include
all node IDs, even those used in previous rounds. All ties are resolved
deterministically.

For all sufficiently large $n$, the cap $M_n$ is not reached, and the
algorithm stops after at most $C_1\sqrt{\log n}$ accepted replacements,
where $C_1$ depends only on the constants in
\eqref{eq:swap-profile-box}. At its endpoint, put
\[
 b_n=a_n+\zeta_n+\frac{\log K}{nP_{\max}}\longrightarrow0.
\]
There is a constant $C_2$, depending only on the constants in the
assumptions and $K$, such that
\begin{align}
 \max_a\min_kD(p_{a*}\Vert\widehat p_{k*})
   &\le C_2b_n nP_{\max},
   \label{eq:swap-population-loss}\\
 J(\widehat p),\ F(\widehat\pi,\widehat p)
   &\le C_2b_n n^2P_{\max}.
   \label{eq:swap-endpoint-cost}
\end{align}
After one common permutation of the fitted components,
\begin{equation}
 \max_a\|\widehat p_{a*}-p_{a*}\|_1
 \le C_2\sqrt{b_n}\,nP_{\max}.
 \label{eq:swap-center-matching}
\end{equation}
For the endpoint's exact E-step responsibilities $R$, the same
permutation gives a deterministic $\gamma_n\to0$ such that
\begin{equation}
 \frac1n\sum_{i,a}|R_{ia}-Z_{ia}|\le\gamma_n.
 \label{eq:swap-weak-responsibilities}
\end{equation}
The statement is uniform over the starting fit, all selected candidates,
and the finite EM stopping rules.
\end{lemma}

\begin{proof}
Write $\mathcal B_n=[cP_{\max},CP_{\max}]^n$ and enlarge $C$ if
necessary without changing $c$. All starting and proposed profiles lie
in $\mathcal B_n$. Each profile M-step is a convex combination of rows
of $q$, and retaining the previous profile in a zero-mass component
preserves this property. Thus the same box contains every profile at
every replacement round. Uniform weights are feasible because
$f<1/K$. Exact EM never decreases $\mathcal L$, including with the
floored simplex, by Proposition~\ref{prop:em-monotonicity}.

\emph{Score identities and bounds uniform over fitted profiles.}
The scalar derivatives of $D(x\Vert y)$ are $\log(x/y)$ and
$1-x/y$. They are uniformly bounded on the common box. Consequently,
for a constant $L$ independent of $n$ and all $x,y,x',y'\in\mathcal B_n$,
\begin{align}
 0\le D(x\Vert y)&\le LnP_{\max},
 \label{eq:swap-bounded-divergence}\\
 |D(x\Vert y)-D(x'\Vert y')|
 &\le L\{\|x-x'\|_1+\|y-y'\|_1\}.
 \label{eq:swap-uniform-lipschitz}
\end{align}
These bounds apply to random fitted profiles and extreme candidate rows;
they make no assumption that a selected candidate is close to its true
class profile. Define the observation-only constant
\[
 S_q=\sum_{i,j}(q_{ij}\log q_{ij}-q_{ij}).
\]
Since $\ell_i(\widehat p_{k*})=\ell_i(q_i)-D(q_i\Vert\widehat p_{k*})$,
the working objective satisfies the exact identities and inequalities
\begin{equation}
 \mathcal L=S_q-F,\qquad
 0\le J\le F,\qquad
 F(\mathbf1/K,\widehat p)\le J(\widehat p)+n\log K.
 \label{eq:swap-saturation-identities}
\end{equation}
For example, the last inequality follows by retaining a minimizing
term in each row's sum. It does not compare weights from different
fits by an unstated monotonicity assumption.

Introduce the population hard cost
\[
 \overline J(\widehat p)=
 \sum_{a=1}^K n_a\min_kD(p_{a*}\Vert\widehat p_{k*}).
\]
For any fixed $K$-tuple $\widehat p\in\mathcal B_n^K$, the elementary
inequality $|\min_k u_k-\min_k v_k|\le\max_k|u_k-v_k|$ and
\eqref{eq:swap-uniform-lipschitz} imply
\begin{equation}
 \sup_{\widehat p\in\mathcal B_n^K}
 |J(\widehat p)-\overline J(\widehat p)|
 \le La_n n^2P_{\max}.
 \label{eq:swap-uniform-population-approximation}
\end{equation}
The supremum is the reason that adaptively selected centers and
data-dependent stopping do not require additional probability bounds.
For every true class $a$, averaging \eqref{eq:swap-average-error} over
its at least $\pi_{\min}n$ nodes supplies a candidate $c(a)$ in that
class with
\begin{equation}
 \|q_{c(a)}-p_{a*}\|_1
 \le (a_n/\pi_{\min})nP_{\max}.
 \label{eq:swap-existing-good-candidate}
\end{equation}
These candidates are used only in the proof; the algorithm does not
know their labels.

\emph{A population replacement supplied by pigeonhole.}
Fix the current $K$ centers. Assign each of the $K$ true profiles to a
minimizing center, using any deterministic tie rule, and write
$k(a)$ for this assignment and
$v_a=D(p_{a*}\Vert\widehat p_{k(a),*})$.
If $k(a)$ is assigned no true profile other than $a$, remove $k(a)$.
If $k(a)$ is assigned at least two true profiles, at most $K-1$ centers
can have any assigned profile, so choose an unassigned center to remove.
In either case, every true class $b\ne a$ retains its assigned minimizing
center. Inserting the exact profile $p_{a*}$ would therefore lower
$\overline J$ by at least $n_av_a$, because the new loss of class $a$
is zero and the losses of the other classes cannot increase. Inserting
the actual row $q_{c(a)}$ instead, \eqref{eq:swap-existing-good-candidate}
and \eqref{eq:swap-uniform-lipschitz} give a replacement configuration
$\widehat p^{\mathrm{trial}}$ satisfying
\begin{equation}
 \overline J(\widehat p^{\mathrm{trial}})
 \le\overline J(\widehat p)-n_av_a
                  +L\pi_{\min}^{-1}a_n n^2P_{\max}.
 \label{eq:swap-pigeonhole-improvement}
\end{equation}
Here the inserted profile may benefit other classes; such benefits are
unneeded for the inequality. Exact score ties cause no difficulty because
the retained minimizing center is still available to every other class.

\emph{The implemented gain finds a replacement at least as good.}
For a removed component $r$, write
\[
 J_{-r}=\sum_i\min_{k\ne r}D(q_i\Vert\widehat p_{k*}),
 \qquad
 h_i^{-r}=\max_{k\ne r}\ell_i(\widehat p_{k*}).
\]
Directly from the score identity,
\begin{equation}
 J(\widehat p^{(r,c)})
 =J_{-r}-\sum_i(\ell_i(q_c)-h_i^{-r})_+,
 \label{eq:swap-global-gain-exact}
\end{equation}
where $\widehat p^{(r,c)}$ replaces $r$ by $q_c$. Thus
\eqref{eq:swap-candidate-rule} is an exact minimization of the
post-replacement hard cost for that deletion. In particular the candidate
chosen for the deletion in \eqref{eq:swap-pigeonhole-improvement} has
empirical hard cost no larger than the proof candidate's cost. Applying
\eqref{eq:swap-uniform-population-approximation} to both configurations
gives, for a constant $L_1$,
\[
 J(\widehat p^{(r,c_r)})
 \le J(\widehat p)-n_av_a+L_1a_n n^2P_{\max}.
\]
Uniform proposal weights and \eqref{eq:swap-saturation-identities},
followed by monotonicity during its fitting, imply that its endpoint has
\begin{equation}
 F_{\mathrm{proposal}}
 \le F_{\mathrm{current}}-n_av_a
                  +L_1a_n n^2P_{\max}+n\log K.
 \label{eq:swap-soft-improvement}
\end{equation}
This inequality holds for every $a$ and its indicated deletion, so the
best evaluated endpoint is at least as good. No purity, convergence,
or basin assumption on the growing initialization has been used.

\emph{The cap is inactive and stopping bounds every population loss.}
The common box gives $cn^2P_{\max}\le T_n\le Cn^2P_{\max}$.
For any feasible weights, \eqref{eq:swap-bounded-divergence} gives
$0\le F\le Ln^2P_{\max}$, because every exponential in each row's
weighted sum is at least $e^{-LnP_{\max}}$. Each accepted replacement
decreases $F$ by at least $\zeta_nT_n$. There are therefore at most
$L/(c\zeta_n)$ accepted rounds. Since this is
$O(\sqrt{\log n})<M_n$ eventually, the prescribed cap is inactive for
all sufficiently large $n$ under the stated deterministic assumptions.
The algorithm must terminate by rejection of all $K+1$ proposals.

At that endpoint, \eqref{eq:swap-soft-improvement} and rejection give
\[
 n_av_a\le L_1a_n n^2P_{\max}+n\log K+\zeta_nT_n.
\]
Dividing by $n_a\ge\pi_{\min}n$ proves
\eqref{eq:swap-population-loss}. Summing over $a$ and using
\eqref{eq:swap-uniform-population-approximation} proves the asserted
bound on $J$. The rejected current-centers uniform-weight proposal has
endpoint cost at most $J+n\log K$, so its rejection also implies
\[
 F_{\mathrm{current}}\le J+n\log K+\zeta_nT_n.
\]
This proves the bound on $F$ in \eqref{eq:swap-endpoint-cost}. The
additional proposal is needed: small current weights can otherwise
make $F$ large even when $J$ and all center errors are small.

\emph{A single center matching and weak responsibilities.}
The Hessian of $\sum_jx_j\log x_j-x_j$ is
$\operatorname{diag}(1/x_j)$. Strong convexity on $\mathcal B_n$, and
then Cauchy--Schwarz, give
\begin{equation}
 D(x\Vert y)\ge\frac{\|x-y\|_2^2}{2CP_{\max}}
              \ge\frac{\|x-y\|_1^2}{2CnP_{\max}}.
 \label{eq:swap-strong-convexity}
\end{equation}
By \eqref{eq:swap-population-loss}, each true class has a fitted
center within $C_3\sqrt{b_n}\,nP_{\max}$ in $\ell_1$.
If two distinct classes chose the same center, their mutual distance
would be at most $2C_3\sqrt{b_n}\,nP_{\max}$, contrary to
\eqref{eq:swap-profile-separation} for large $n$. With exactly $K$
true profiles and $K$ fitted profiles, this injection is a bijection.
It supplies the common permutation in \eqref{eq:swap-center-matching}.

Set $t_n=\sqrt{a_n}+(nP_{\max})^{-1}$ and call row $i$ typical when
$\|q_i-p_{z_i,*}\|_1\le t_nnP_{\max}$. By
\eqref{eq:swap-average-error}, the fraction of nontypical rows is at
most $a_n/t_n\to0$. For a typical row in class $a$ and any incorrectly
matched component $k\ne a$, separation and the center matching imply
\[
 \|q_i-\widehat p_{k*}\|_1\ge (s/2)nP_{\max}
\]
for all sufficiently large $n$. By \eqref{eq:swap-strong-convexity},
its deviance is at least $c_*nP_{\max}$, with
$c_*=s^2/(8C)>0$. The exact E-step satisfies
\[
 R_{ik}=\frac{\widehat\pi_ke^{-D(q_i\Vert\widehat p_{k*})}}
               {\sum_l\widehat\pi_le^{-D(q_i\Vert\widehat p_{l*})}}.
\]
Its row-wise variational identity is
\[
 -\log\sum_k\widehat\pi_ke^{-D(q_i\Vert\widehat p_{k*})}
 =\sum_kR_{ik}D(q_i\Vert\widehat p_{k*})
       +\operatorname{KL}(R_i\Vert\widehat\pi)
 \ge\sum_kR_{ik}D(q_i\Vert\widehat p_{k*}).
\]
Summing over typical rows and using
$\sum_k|R_{ik}-Z_{ik}|=2\sum_{k\ne z_i}R_{ik}\le2$ gives
\[
 \frac1n\sum_{i,k}|R_{ik}-Z_{ik}|
 \le\frac{2a_n}{t_n}+
       \frac{2F}{c_*n^2P_{\max}}
 \le\frac{2a_n}{t_n}+\frac{2C_2b_n}{c_*}
 =:\gamma_n\longrightarrow0.
\]
The identities remain valid for zero weights with the usual convention
that the corresponding responsibilities vanish. In particular, no
unstated lower bound on the endpoint's fitted mixture weights is
required for this weak-start argument. All bounds are simultaneous for
every configuration in the box, completing the deterministic proof.
\end{proof}

\paragraph{Proof of Theorem~\ref{thm:certified-likelihood}.}
Fix a confidence constant $H$ sufficiently large that the exceptional
probability in Lemma~\ref{lem:full-spectral} is negligible relative to
$e^{-I_n}$. Under the density-uniform form of that lemma, its probability
is $O(e^{-HnP_{\max}})$, and the upper bound
$I_n\le C_I nP_{\max}$ allows any fixed $H>C_I$.
Work on its spectral-profile event. The positive probability ratios,
positive limiting community proportions, and the spectral conditioning
assumption give \eqref{eq:swap-profile-box} and
\eqref{eq:swap-profile-separation}, with constants that may depend on
$H$. Its average-profile bound is exactly
\eqref{eq:swap-average-error}. The growing fit starts at a row mean and
only appends rows and takes convex profile updates, so all of its
endpoint profiles lie in the required row convex hull.

Apply Lemma~\ref{lem:iterative-swap-certification} to the deterministic
replacement stage of Algorithm~1. For every sufficiently large $n$ on
this event, its round cap is inactive and its endpoint E-step obeys
\eqref{eq:swap-weak-responsibilities}. Thus the local initialization
condition \eqref{eq:local-initialization} holds with a deterministic
$\gamma_n\to0$ and $\delta_n=O(e^{-HnP_{\max}})$.
There is no additional randomized-initialization failure. The mandatory
final profile and weight M-step and
Theorem~\ref{thm:local-spectral-em} give
\[
 \E\Mis(\widehat z,z)\le e^{-(1+o(1))I_n}.
\]
The same theorem covers any further exact updates and their stated
finite stopping rules. If $\liminf I_n/\log n>1$, multiplying the
expected error fraction by $n$ and applying Markov's inequality proves
exact recovery with probability tending to one. All replacement
proposals use only $q$, hence only the retained eigenpairs, and the
number of replacement rounds and the fitting iteration budgets are
polynomial by construction. No recovery property of the unrepaired
growing endpoint is needed. \qed
